\documentclass[11pt,a4paper]{ctexart}
\usepackage[a4paper,margin=1.8cm,headheight=14pt]{geometry}
\usepackage{amsmath,amssymb,mathtools}
\usepackage{booktabs,tabularx,longtable,array,multirow}
\usepackage{enumitem}
\usepackage{tikz}
\usetikzlibrary{arrows.meta,positioning,shapes.geometric,fit,calc}
\usepackage[most]{tcolorbox}
\usepackage{xcolor}
\usepackage{fancyhdr}
\usepackage{hyperref}
\usepackage{microtype}
\setlength{\parindent}{2em}
\setlength{\parskip}{0.35em}
\setlist[itemize]{itemsep=0.2em,topsep=0.3em}
\setlist[enumerate]{itemsep=0.2em,topsep=0.3em}
\hypersetup{colorlinks=true,linkcolor=blue!45!black,urlcolor=blue!50!black}

\definecolor{mainblue}{RGB}{45,86,140}
\definecolor{deepblue}{RGB}{30,62,102}
\definecolor{softblue}{RGB}{238,245,252}
\definecolor{softgreen}{RGB}{238,249,243}
\definecolor{softorange}{RGB}{253,246,233}
\definecolor{softred}{RGB}{253,239,239}
\definecolor{softgray}{RGB}{246,247,249}
\definecolor{deepgreen}{RGB}{35,105,75}
\definecolor{deeporange}{RGB}{165,98,22}
\definecolor{deepred}{RGB}{150,55,55}

\pagestyle{fancy}
\fancyhf{}
\lhead{408 操作系统方法论}
\rhead{OS 科内桥梁与跨科接口手册}
\cfoot{\thepage}
\renewcommand{\headrulewidth}{0.4pt}

\newtcolorbox{corebox}[1][]{colback=softblue,colframe=mainblue,title=#1,fonttitle=\bfseries,boxrule=0.8pt,arc=2mm}
\newtcolorbox{exam}[1][]{colback=softorange,colframe=deeporange,title=#1,fonttitle=\bfseries,boxrule=0.8pt,arc=2mm}
\newtcolorbox{warn}[1][]{colback=softred,colframe=deepred,title=#1,fonttitle=\bfseries,boxrule=0.8pt,arc=2mm}
\newtcolorbox{eng}[1][]{colback=softgreen,colframe=deepgreen,title=#1,fonttitle=\bfseries,boxrule=0.8pt,arc=2mm}
\newtcolorbox{method}[1][]{colback=softgray,colframe=black!45,title=#1,fonttitle=\bfseries,boxrule=0.6pt,arc=2mm}

\newcommand{\kw}[1]{\textbf{#1}}
\newcommand{\code}[1]{\texttt{#1}}

\title{\vspace{-1.2cm}\textbf{操作系统科内桥梁与跨科接口}\\[0.4em]
\large 破除模块孤岛：跨 Topic 机制冲突、控制权交接与跨科连接}
\author{408 操作系统专题手册 OS-Bridge}
\date{}

\begin{document}
\maketitle

\begin{corebox}[母问题：为什么单独掌握每个 Topic 依然做不对综合大题？]
在真实的操作系统和 408 综合大题中，机制绝不是独立运行的。当调度遇上并发锁、当 CPU 试图访问 I/O 设备、当进程复制遇到文件描述符，不同子系统之间的\kw{规则冲突与控制权交接}便构成了系统设计的核心挑战。

本手册专为连接各个模块而建立：
\[
\boxed{
\text{Process / Scheduling}
\xleftrightarrow{\quad\text{Bridges}\quad}
\text{Concurrency / VM / I/O / File System / CPU Hardware}
}
\]
\end{corebox}

\tableofcontents
\newpage

\section{Bridge 1：Priority Inversion（优先级倒置）——调度与并发锁的冲突}

\subsection{冲突的产生：当 Scheduling 遇上 Mutual Exclusion}

操作系统有两个基本机制：
\begin{itemize}
  \item \kw{进程调度器 (Scheduling)} 规定：优先级高的任务应优先获得 CPU（$P_{High} > P_{Med} > P_{Low}$）；
  \item \kw{并发互斥锁 (Lock / Mutex)} 规定：若锁已被持用，等待该锁的任何任务必须阻塞，直到锁被释放。
\end{itemize}

当高优先级任务 $P_{High}$ 试图获取被低优先级任务 $P_{Low}$ 持有的锁时，$P_{High}$ 被迫进入阻塞等待（符合互斥逻辑）。

但在此时，若中优先级任务 $P_{Med}$（不需要该锁）变为可运行状态，调度器会选择抢占 $P_{Low}$ 的 CPU！

这引发了极其严重的系统逆转：
\[
\boxed{
P_{High} \quad\text{间接等待}\quad P_{Med}
}
\]
由于 $P_{Low}$ 无法获得 CPU 运行，它无法释放锁，导致最高优先级的 $P_{High}$ 无法运行！这就是\kw{优先级倒置 (Priority Inversion)}。

\subsection{优先级继承 (Priority Inheritance) 解决方案}

为了打破这一僵局，系统引入\kw{优先级继承}协议：
\[
\boxed{
\text{当 } P_{High} \text{ 阻塞等待 } P_{Low} \text{ 持有的锁时，} P_{Low} \text{ 临时继承 } P_{High} \text{ 的高优先级}
}
\]

\begin{center}
\begingroup
\small
\renewcommand{\arraystretch}{1.25}
\setlength{\tabcolsep}{6pt}
\begin{tabularx}{\linewidth}{>{\bfseries}p{0.22\linewidth} p{0.36\linewidth} p{0.36\linewidth}}
\toprule
解决机制 & 核心工作原理 & 评价与工程代价\\
\midrule
Priority Inheritance (优先级继承) & 锁持有人 $P_{Low}$ 临时升至等待者最高优先级 $P_{High}$，防止中优先级 $P_{Med}$ 抢占；释放锁后恢复原优先级 & 动态修改任务优先级，开销较低；但不能完全消除死锁风险\\
Priority Ceiling (优先级天花板) & 给锁设定最高可能访问者的优先级，一旦获取锁即提升至天花板优先级 & 进一步预防死锁与多次阻塞，适用于硬实时系统\\
\bottomrule
\end{tabularx}
\endgroup
\end{center}

\section{Bridge 2：MMIO vs Port I/O——CPU $\times$ OS $\times$ I/O 的跨科连接}

\subsection{CPU 怎样看到外部设备控制器}

在 408 计算机组成原理与操作系统交汇处，核心问题是：\kw{CPU 指令如何对设备控制器中的控制/状态/数据寄存器进行读写？}

存在两种主要的编址方式：

\begingroup
\small
\renewcommand{\arraystretch}{1.25}
\setlength{\tabcolsep}{6pt}
\begin{tabularx}{\linewidth}{>{\bfseries}p{0.25\linewidth} p{0.35\linewidth} p{0.35\linewidth}}
\toprule
编址方式 & 核心实现原理 & 优缺点与指令差异\\
\midrule
Memory-Mapped I/O (MMIO 内存映射 I/O) & 将设备寄存器映射到常规物理内存地址空间；CPU 使用统一的 \code{load/store} 访存指令访问设备 & 统一指令集，可利用 MMU 页表进行权限保护；但占用物理地址空间，需防止 CPU/Cache 乱序优化（需设为 Non-cacheable）\\
Isolated / Port-Mapped I/O (独立/端口 I/O) & 设备使用独立的 64KB I/O 端口地址空间；CPU 使用专门的 I/O 指令（如 x86 \code{in/out}） & 内存与 I/O 隔离；但依赖专有指令，无法直接享用 MMU 高级保护与灵活寻址\\
\bottomrule
\end{tabularx}
\endgroup

\section{Bridge 3：Process $\times$ File——引用交接与分层模型}

当进程执行 \code{fork()}、\code{dup()} 或 \code{close()} 时，文件控制权在三级结构间交接：

\[
\boxed{
\text{Task / PCB (Process-private Table)}
\xrightarrow{\text{File Descriptor } fd}
\text{Open File Description (OFD / System-wide Table)}
\xrightarrow{\text{Inode Index}}
\text{Inode Object (FileSystem)}
}
\]

\begin{itemize}
  \item \kw{fork()}：子进程复制父进程的 \code{fd table}，使得父子进程的相同 $fd$ 指向同一个 \code{OFD}（共享文件偏移量 \code{offset}）；
  \item \kw{dup()}：在同一进程内创建新 $fd$，指向相同的 \code{OFD}（共享 \code{offset}）；
  \item \kw{open()} 同一文件两次：创建两个独立 \code{OFD}（拥有各自独立的 \code{offset}），但底层指向同一个 \code{inode}。
\end{itemize}

\section{Bridge 4：VM $\times$ File——mmap 与 Page Cache 的结合}

\code{mmap()} 建立了虚拟内存与文件系统的直接接口：

\[
\boxed{
\text{VMA (Virtual Memory Area)}
\xrightarrow{\text{Page Fault}}
\text{Page Cache (RAM Physical Frame)}
\xrightarrow{\text{Dirty Writeback}}
\text{Inode Block Mapping (Disk)}
}
\]

通过把文件的 \code{Page Cache} 物理页直接映射到进程的虚拟页表中，避免了传统 \code{read()} / \code{write()} 在内核缓冲区与用户缓冲区之间的显式数据复制！

\end{document}
